% Unidad U006D. Registro dodecafásico integral y descomposición 1+5+6.

\chapter{Registro dodecafásico integral, extensión cíclica y descomposición \(1+5+6\)}
\label{chap:registro-dodecafasico-integral}

El periodo observable de ciento ocho posiciones contiene doce sectores de
nueve posiciones. Esta igualdad de cardinales no basta para describir el
estado temporal: deben distinguirse el soporte ordenado, la ley cíclica, el
índice sectorial y el registro enriquecido de transportes, orientaciones,
acarreos e incidencias. Esta unidad construye los cuatro objetos y demuestra
la reconstrucción integral de las doce coordenadas de memoria.

\section{Soporte ordenado y coordenada cíclica}

Para \(m\in\{1,\ldots,12\}\), definimos
\begin{equation}
 E_m:=\{9(m-1)+1,\ldots,9m\}.
 \label{eq:u006d-bloques-nonadicos}
\end{equation}
El soporte ordenado y el grupo cíclico son
\begin{equation}
 \Gamma_{108}:=
 E_1\sqcup\cdots\sqcup E_{12},
 \qquad
 \Theta_{108}:=\mathbb Z/108\mathbb Z.
 \label{eq:u006d-gamma-theta}
\end{equation}
\(\Gamma_{108}\) conserva posición, orden y pertenencia a bloque;
\(\Theta_{108}\) conserva la ley de traslación. No son el mismo objeto.

Para \(\theta\in\Theta_{108}\), sea
\(\widetilde\theta\in\{0,\ldots,107\}\) su representante. Definimos
\begin{equation}
 \beta(\theta):=
 \left[
 \left\lfloor\frac{\widetilde\theta}{9}\right\rfloor
 \right]_{12},
 \qquad
 r(\theta):=1+(\widetilde\theta\bmod9).
 \label{eq:u006d-coordenadas-bloque-residuo}
\end{equation}
La primera coordenada localiza el sector dodecafásico; la segunda publica
la posición positiva \(1,\ldots,9\) dentro del sector.

\section{Extensión cíclica no escindida}

Sean \(C_n=\mathbb Z/n\mathbb Z\). Las aplicaciones
\begin{equation}
 \iota:C_{12}\longrightarrow C_{108},
 \quad \iota([b]_{12})=[9b]_{108},
 \qquad
 p:C_{108}\longrightarrow C_9,
 \quad p([\theta]_{108})=[\theta]_9
 \label{eq:u006d-mapas-extension}
\end{equation}
forman la sucesión
\begin{equation}
 0\longrightarrow C_{12}
 \xrightarrow{\;\iota\;}C_{108}
 \xrightarrow{\;p\;}C_9
 \longrightarrow0.
 \label{eq:u006d-extension}
\end{equation}
Para la sección conjuntista
\[
 s([r]_9):=[\widetilde r]_{108},
 \qquad 0\leq\widetilde r<9,
\]
el defecto de aditividad es
\begin{equation}
 c(r,r'):=
 \left[
 \left\lfloor
 \frac{\widetilde r+\widetilde r'}9
 \right\rfloor
 \right]_{12},
 \qquad
 s(r)+s(r')
 =s(r+r')+\iota(c(r,r')).
 \label{eq:u006d-cociclo-extension}
\end{equation}

\begin{theorem}[Exactitud, no escisión y clase cohomológica]
\label{thm:u006d-extension-no-escindida}
La sucesión \eqref{eq:u006d-extension} es exacta y no se escinde. Para la
acción trivial,
\begin{equation}
 H^2(C_9,C_{12})
 \simeq C_{12}/9C_{12}
 \simeq C_3,
 \label{eq:u006d-cohomologia-extension}
\end{equation}
y el cociclo \eqref{eq:u006d-cociclo-extension} representa su clase
generadora.
\end{theorem}

\begin{proof}
La imagen de \(\iota\) es el subgrupo de múltiplos de nueve, que coincide
con \(\ker p\); \(p\) es sobreyectiva. Esto prueba la exactitud. Si la
extensión se escindiera, \(C_{108}\) sería isomorfo a
\(C_{12}\times C_9\), pero
\[
 \exp(C_{108})=108,
 \qquad
 \exp(C_{12}\times C_9)=\operatorname{mcm}(12,9)=36.
\]
La identidad de cociclo es la división euclídea de
\(\widetilde r+\widetilde r'\) por nueve. Para un grupo cíclico con acción
trivial,
\(H^2(C_9,C_{12})\simeq C_{12}/9C_{12}\); el acarreo unitario proyecta a
\([1]\), que genera el cociente de orden tres.
\end{proof}

La no escisión tipa el retorno:
\[
 g_0(\theta+9)=g_0(\theta),
 \qquad
 \beta(\theta+9)=\beta(\theta)+1\pmod{12}.
\]
Nueve iteraciones cierran la fase visible y trasladan un sector; ciento
ocho cierran las dos coordenadas del reloj. El estado enriquecido conserva
además la memoria vertical, el cilindro, la frontera y la genealogía.

\section{Registro temporal enriquecido}

\begin{definition}[Registro dodecafásico orientado]
\label{def:u006d-registro-dodecafasico}
Para una historia enriquecida, sea \(\tau_m\) la subruta ordenada sobre
\(E_m\), \(\sigma_m\) su orientación, \(\kappa_m\) el acarreo que cruza la
frontera del bloque y \(\Lambda_m\) su libro local de incidencias. El
registro completo es
\begin{equation}
 \mathcal R_{12}:=
 \bigl(
 (E_m,\tau_m,\sigma_m,\kappa_m,\Lambda_m)
 \bigr)_{m=1}^{12}.
 \label{eq:u006d-registro-r12}
\end{equation}
\end{definition}

Se distinguen así cuatro niveles:
\begin{equation}
 \boxed{
 \mathcal R_{12}
 \neq C_{12}
 \neq\Theta_{108}
 \neq\Gamma_{108}.}
 \label{eq:u006d-cuatro-objetos}
\end{equation}
El primero es un registro enriquecido; el segundo, un índice de sector; el
tercero, un grupo de traslaciones; el cuarto, un soporte ordenado. Existen
proyecciones entre ellos, pero no identificaciones.

Sea \(q_{\rm mem}\in\mathbb Z_{\geq0}\) el número no acotado de vueltas
nonádicas. La coordenada dodecafásica es sólo su residuo
\[
 \beta=[q_{\rm mem}]_{12}.
\]
Tras ciento ocho pasos,
\[
 \beta\longmapsto\beta,
 \qquad
 q_{\rm mem}\longmapsto q_{\rm mem}+12.
\]
El calendario retorna; la memoria vertical no se reinicia.

\section{Extractor firmado de eventos}

Sea \((d_t)_{t=1}^{108}\) el registro digital sobre
\(\Gamma_{108}\). Definimos
\begin{equation}
 \mathcal D_{\rm evt}:=\{2,4,6,8\},
 \qquad
 \Sigma_{12}:=(+,+,+,-,-,-,+,+,+,-,-,-).
 \label{eq:u006d-datos-eventos}
\end{equation}
Para \(m=0,\ldots,11\), sea
\begin{equation}
 e_m:=
 \Sigma_{12}(m+1)\,
 \#\{t\in E_{m+1}:d_t\in\mathcal D_{\rm evt}\}
 \in\mathbb Z.
 \label{eq:u006d-eventos-firmados}
\end{equation}
La aplicación conserva el censo local y la orientación sectorial como
factores separados.

Emparejamos las dos semicoronas mediante
\begin{equation}
 p_i:=e_i+e_{i+6},
 \qquad
 a_i:=e_i-e_{i+6},
 \qquad 0\leq i\leq5.
 \label{eq:u006d-pares-antipodales}
\end{equation}
Definimos el censo total, las cinco diferencias respecto del sexto par y
las seis antisimetrías:
\begin{equation}
 s_C:=\sum_{i=0}^{5}p_i,
 \qquad
 M_i:=p_i-p_5\quad(0\leq i\leq4),
 \qquad
 A_6:=(a_0,\ldots,a_5).
 \label{eq:u006d-coordenadas-uno-cinco-seis}
\end{equation}

\begin{definition}[Lector integral \(1+5+6\)]
\label{def:u006d-lector-integral}
El lector lineal es
\begin{equation}
 \mathcal L_{12}:\mathbb Z^{12}\longrightarrow\mathbb Z^{12},
 \qquad
 \mathcal L_{12}(e)
 :=
 (s_C,M_0,\ldots,M_4,a_0,\ldots,a_5).
 \label{eq:u006d-lector-l12}
\end{equation}
La distribución de coordenadas es
\begin{equation}
 \underbrace{1}_{s_C}
 +\underbrace{5}_{M_0,\ldots,M_4}
 +\underbrace{6}_{a_0,\ldots,a_5}
 =12.
 \label{eq:u006d-uno-cinco-seis}
\end{equation}
\end{definition}

\begin{theorem}[Reconstrucción exacta e imagen integral]
\label{thm:u006d-inversa-l12}
El lector \(\mathcal L_{12}\) es inyectivo. Sobre su imagen,
\begin{align}
 p_5&=
 \frac{s_C-\sum_{i=0}^{4}M_i}{6},
 &
 p_i&=p_5+M_i\quad(0\leq i\leq4),
 \label{eq:u006d-inversa-p}\\
 e_i&=\frac{p_i+a_i}{2},
 &
 e_{i+6}&=\frac{p_i-a_i}{2}
 \quad(0\leq i\leq5).
 \label{eq:u006d-inversa-e}
\end{align}
Una tupla de salida pertenece a \(\operatorname{im}\mathcal L_{12}\) si y
sólo si
\begin{equation}
 s_C-\sum_{i=0}^{4}M_i\equiv0\pmod6
 \label{eq:u006d-condicion-divisibilidad}
\end{equation}
y, para los \(p_i\) reconstruidos,
\begin{equation}
 p_i\equiv a_i\pmod2
 \qquad(0\leq i\leq5).
 \label{eq:u006d-condicion-paridad}
\end{equation}
\end{theorem}

\begin{proof}
De \(M_i=p_i-p_5\) se obtiene
\[
 s_C=\sum_{i=0}^{4}(p_5+M_i)+p_5
 =6p_5+\sum_{i=0}^{4}M_i,
\]
que da la primera ecuación de \eqref{eq:u006d-inversa-p}; las demás son
inmediatas. Sumar y restar
\(p_i=e_i+e_{i+6}\) y \(a_i=e_i-e_{i+6}\) produce
\eqref{eq:u006d-inversa-e}. La divisibilidad por seis y las seis
congruencias de paridad son precisamente las condiciones necesarias y
suficientes para que todas las coordenadas reconstruidas sean enteras.
\end{proof}

La división por seis aparece después del censo orientado, no durante él. Las
condiciones integrales de esta unidad son exactamente
\eqref{eq:u006d-condicion-divisibilidad} y
\eqref{eq:u006d-condicion-paridad}; no se atribuye aquí una forma normal de
Smith a una matriz que no haya sido previamente definida.

\section{Dos palabras dodecafásicas de distinta procedencia}

Debe conservarse la distinción
\begin{equation}
 U_{12}^{\rm cat}:=U_6\Vert U_6,
 \qquad
 U_{12}^{\rm sgn}:=\mathcal H_{12}(K).
 \label{eq:u006d-dos-palabras-doce}
\end{equation}
La primera concatena dos emisiones catalogales de longitud seis. La segunda
es una coordenada firmada ligada reversiblemente a \(K\) por el operador
\(\mathcal H_{12}\). Su longitud común no implica una misma procedencia,
semántica o ley de transporte.

\begin{criterion}[Criterios de refutación]
\label{crit:u006d-refutacion}
La construcción queda refutada si la sucesión
\eqref{eq:u006d-extension} no es exacta, si su clase cohomológica es
trivial, si una tupla de la imagen viola
\eqref{eq:u006d-condicion-divisibilidad} o
\eqref{eq:u006d-condicion-paridad}, si las fórmulas inversas no recuperan
los doce enteros iniciales, o si el retorno de \(\beta\) se interpreta como
reinicio de \(q_{\rm mem}\).
\end{criterion}
